\documentclass[conference]{IEEEtran}

\usepackage{booktabs}
\usepackage{enumitem}
\usepackage{fvextra}
\usepackage{hyperref}
\usepackage{microtype}
\usepackage{tabularx}
\usepackage{tagging}
\usepackage{sc26repro}

\hypersetup{colorlinks=true,linkcolor=black,urlcolor=blue}
\setlist{nosep,leftmargin=*}

\begin{document}

\twocolumn[%
{\begin{center}
\Huge
Appendix: Artifact Evaluation
\end{center}}
]

\small

\appendixAE
\enlargethispage{8\baselineskip}

Paper ID: \texttt{pap142}

This workflow evaluates \texttt{A1} (TempGNN and three U280 accelerator
baselines) and \texttt{A2} (\path{results/result.csv} and the paper-figure
generator).

\noindent Artifact:
\href{https://doi.org/10.5281/zenodo.21417187}{10.5281/zenodo.21417187};
source:
\href{https://github.com/TempGNN-Program/TempGNN}{TempGNN-Program/TempGNN};
tag: \texttt{sc26-ae-pap142-v1.0}.

\section{Download and Setup}

\begin{Verbatim}[fontsize=\small,breaklines=true,breakanywhere=true]
git clone https://github.com/TempGNN-Program/TempGNN.git
cd TempGNN
python3 -m venv .venv
source .venv/bin/activate
pip install -r requirements.txt
\end{Verbatim}

\begin{Verbatim}[fontsize=\small,breaklines=true,breakanywhere=true]
Board: AMD/Xilinx Alveo U280, xcu280-fsvh2892-2L-e
Platform: xilinx_u280_gen3x16_xdma_1_202211_1
OS: Ubuntu 22.04 x86_64
XRT: 2.16.204
Vitis/Vivado: 2023.2 (xclbin rebuild only)
GCC: 11.4.0
GNU Make: 4.3
Python: 3.10 or newer
\end{Verbatim}

The U280 account and login instructions are delivered through the private
SC26 submission channel. No credentials are stored in the artifact.

\section{Review Plan}

\begin{table}[h]
\caption{Expected direct review time.}
\centering
\footnotesize
\begin{tabularx}{\columnwidth}{@{}lrrrX@{}}
\toprule
Artifact & Setup & Run & Analysis & Expected result \\
\midrule
\texttt{A1} & 10--20 & 30--90 & $<$5 & Four fresh latency results; TempGNN
lowest at about 1.3 ms \\
\texttt{A2} & $<$5 & $<$1 & 2--5 & Nine CSV/SVG pairs from
\path{results/result.csv} \\
\bottomrule
\end{tabularx}
\end{table}

All times are minutes. Optional xclbin rebuilding is excluded from the direct
review path.

\section{Execute A1: U280 Latency}

\begin{Verbatim}[fontsize=\small,breaklines=true,breakanywhere=true]
source /opt/xilinx/xrt/setup.sh
xrt-smi examine
make ae-core-u280 U280_CORE_DEVICE=0 U280_CORE_REPETITIONS=3
\end{Verbatim}

\noindent The command performs:
\begin{enumerate}
\item \texttt{A1-T1}: check the four packaged U280 artifacts;
\item \texttt{A1-T2}: load six datasets and four models;
\item \texttt{A1-T3}: run TempGNN, MATG, ViTeGNN, and RTGA three times;
\item \texttt{A1-T4}: write per-implementation raw measurement CSVs; and
\item \texttt{A1-T5}: aggregate latency and update the AE report.
\end{enumerate}

\begin{Verbatim}[fontsize=\small,breaklines=true,breakanywhere=true]
Datasets: WK, MC, RT, LM, WT, GT
Models: JODIE, TGN, TGAT, APAN
Repetitions: 3
\end{Verbatim}

Fresh results are written to
\path{results/reviewer_u280_runs/<run-id>/}.

\section{Analyze A1}

The run is successful when:
\begin{enumerate}
\item the command exits with status zero;
\item fresh rows exist for all four implementations;
\item all six datasets and four models are present; and
\item the aggregate latency follows the expected U280 behavior in
Table~\ref{tab:ae-latency}.
\end{enumerate}

\begin{table}[h]
\caption{Expected U280 mean latency.}
\label{tab:ae-latency}
\centering
\begin{tabular}{lr}
\toprule
Implementation & Latency (ms) \\
\midrule
TempGNN & 1.296553 \\
MATG & 9.966618 \\
ViTeGNN & 2.869752 \\
RTGA & 4.192824 \\
\bottomrule
\end{tabular}
\end{table}

The principal behavior is that TempGNN has the lowest mean latency and remains
close to 1.3 ms. This supports the target-centric execution, DDTC/OATS
mechanisms, and U280 accelerator comparison in \texttt{C1-C4}.

\section{Execute and Analyze A2}

\begin{Verbatim}[fontsize=\small,breaklines=true,breakanywhere=true]
python3 -m scripts.reproduce_paper_figures
\end{Verbatim}

\noindent The command performs:
\begin{enumerate}
\item \texttt{A2-T1}: read \path{results/result.csv};
\item \texttt{A2-T2}: validate and group rows by paper figure;
\item \texttt{A2-T3}: generate a CSV and SVG for each figure; and
\item \texttt{A2-T4}: write \path{figure_data_manifest.csv}.
\end{enumerate}

Expected output under \path{results/paper_reproduction/}:
\begin{itemize}
\item Fig.2, Fig.4(a), and Fig.9(b);
\item Fig.10, Fig.11, Fig.12, and Fig.13; and
\item Fig.14(a), Fig.14(b), and the figure-data manifest.
\end{itemize}

The evaluation is successful when all nine CSV/SVG pairs exist and the
manifest records \path{results/result.csv} as their input. These outputs
provide the paper-result views associated with \texttt{C1-C4}.

\section{Baseline Status}

\begin{itemize}
\item MATG: independently reproduced from the IPDPS 2022 paper and its partial
public HLS headers; includes a distinct source, runner, and U280 xclbin.
\item ViTeGNN: independently reproduced from the TPDS 2025 paper; includes a
distinct source, runner, and U280 xclbin.
\item RTGA: independently reproduced from the DAC 2024 paper; includes a
distinct source, runner, and U280 xclbin.
\item Cascade and TGLite-CPU: retained only as packaged paper comparison
records; no fresh execution of those external stacks is claimed.
\item TempGNN: includes HLS kernels, testbenches, Vitis build scripts, XRT host
code, U280 board logs, and its own U280 xclbin.
\end{itemize}

The three baselines are independent, paper-based forward-path reproductions,
not the authors' complete original stacks.

\section{Metric Meaning}

\texttt{DDTC} is dependency-driven TDP construction. \texttt{OATS} is
overlap-aware TDP synchronization.

\texttt{packet\_reuse\_factor} measures how many per-target packet
materializations collapse into unique PHLE packets. \texttt{memory\_bytes} is
packet state/metadata traffic. \texttt{cycles} is the kernel's cycle proxy used
by C-sim and exported fixture stats. HLS reports provide C/RTL latency after
synthesis/cosim.

\section{License}

This artifact is licensed under the Apache License 2.0. See \path{LICENSE}.

\end{document}
